\documentclass[11pt,a4paper,danish]{article}
\usepackage[utf8]{inputenc}
\usepackage[danish]{babel} % danske overskrifter
\usepackage[T1]{fontenc} % fonte (output)
\usepackage{graphicx} % indsættelse af billeder
\usepackage{caption}
\usepackage{subcaption}
\usepackage{wrapfig}
\usepackage{geometry}
\usepackage{titlesec}
\usepackage{fancyhdr}
\usepackage{lastpage}
%\usepackage{biblatex}
\usepackage{hyperref}
\usepackage{gensymb}

\usepackage{setspace}
\usepackage{tocloft}
\usepackage{parskip}

\usepackage[nottoc]{tocbibind}

%\usepackage{subfig}
%\usepackage{inconsolata}
%\usepackage{enumitem}

\titleformat{\paragraph}[runin]
{\bfseries\scshape}{\theparagraph}{1em}{}

% BEGIN JAVA CODE SECTION
\usepackage{listings}
\usepackage{color}

\definecolor{dkgreen}{rgb}{0,0.6,0}
\definecolor{gray}{rgb}{0.5,0.5,0.5}
\definecolor{mauve}{rgb}{0.58,0,0.82}

\lstset{frame=none,
  language=Java,
  aboveskip=3mm,
  belowskip=3mm,
  showstringspaces=false,
  columns=flexible,
  basicstyle={\small\ttfamily},
  numbers=none,
  numberstyle=\tiny\color{gray},
  keywordstyle=\color{blue},
  commentstyle=\color{dkgreen},
  stringstyle=\color{mauve},
  breaklines=true,
  breakatwhitespace=true,
  tabsize=2
}
% END JAVA CODE SECTION

% BEGIN C++ code
\definecolor{listinggray}{gray}{0.9}
\definecolor{lbcolor}{rgb}{0.9,0.9,0.9}
\lstset{
backgroundcolor=\color{lbcolor},
    tabsize=4,
%   rulecolor=,
    language=[GNU]C++,
        basicstyle=\scriptsize,
        upquote=true,
        aboveskip={1.5\baselineskip},
        columns=fixed,
        showstringspaces=false,
        extendedchars=false,
        breaklines=true,
        prebreak = \raisebox{0ex}[0ex][0ex]{\ensuremath{\hookleftarrow}},
        frame=single,
        numbers=left,
        showtabs=false,
        showspaces=false,
        showstringspaces=false,
        identifierstyle=\ttfamily,
        keywordstyle=\color[rgb]{0,0,1},
        commentstyle=\color[rgb]{0.026,0.112,0.095},
        stringstyle=\color[rgb]{0.627,0.126,0.941},
        numberstyle=\color[rgb]{0.205, 0.142, 0.73},
%        \lstdefinestyle{C++}{language=C++,style=numbers}’.
}
\lstset{
    backgroundcolor=\color{lbcolor},
    tabsize=4,
  language=C++,
  captionpos=b,
  tabsize=3,
  frame=lines,
  numbers=left,
  numberstyle=\tiny,
  numbersep=5pt,
  breaklines=true,
  showstringspaces=false,
  basicstyle=\footnotesize,
%  identifierstyle=\color{magenta},
  keywordstyle=\color[rgb]{0,0,1},
  commentstyle=\color{Darkgreen},
  stringstyle=\color{red}
  }
%END C++ code



%\setcounter{secnumdepth}{3}

\geometry{
 a4paper,
 total={210mm,297mm},
 left=20mm,
 right=20mm,
 top=40mm,
 bottom=20mm,
 }

\pagestyle{fancy}
\fancyhf{}
\rhead{Delta, \date{\today}}
\lhead{Rudy Alex Kohn}
\chead{Uge 2}
\cfoot{Side \thepage \hspace{1pt} af \pageref{LastPage}}

% TOC som links
\usepackage{hyperref}
\hypersetup{
    colorlinks,
    citecolor=black,
    filecolor=black,
    linkcolor=black,
    urlcolor=blue
}

\title{Stripetykkelsesmåler\\Afsluttende Rapport For Praktikophold}
\author{Rudy Alex Kohn}
\date{8. Juni 2017}

\setlength{\parindent}{0pt}

\begin{document}

\maketitle
\renewcommand{\contentsname}{Indeks}
\renewcommand{\listfigurename}{Figurliste}
\renewcommand{\figurename}{Figur}
\renewcommand\refname{Referencer}
%PWNED! :> nice!

\begin{center}
Rudy Alex Kohn\\
ruak@force.dk\\
s133235@student.dtu.dk
\end{center}

\vspace{25mm}

\begin{center}
\includegraphics[scale=0.3]{Billeder/kunlogo.png}
\includegraphics[scale=0.15]{Billeder/DELTA_1024px.png}
\end{center}

\newpage
\setlength\cftparskip{-2pt}
%\setlength\cftbeforechapskip{0.0pt}\singlespacing
\tableofcontents
%\singlespacing

\listoffigures
\listoftables

%1:  produktspecifikation. forudsætninger
%Performance  tykkelsesmålinh præsisosn +/- 0.1 mm
%Product cist max 25,000,- kr gerne mindre.
%2 teknologier og teknologi valg – begrundelse
%3 forsøgsopstilling – gennemgang.
%4 databehandling  procesdiagram.
%5 resulateter  - hvordan ser det ud.
%6 Videre arbejde- plan for den sidste del af projektet- hvor langt når vi, dokumentation etc.

\section{Introduktion}

Denne rapport afdækker praktikforløbet hos Delta (Force, Hørsholm). Den beskriver de forskellige faser i projektet, samt belyser de problematikker der opstod i forbindelse med at løse opgaven.
Praktikopholdet hos Delta er baseret på et afklaringsprojekt med henblik på at få afdækket om hvorhvidt det er muligt at kunne måle højden af en vejstripe (termoplast). I den forbindelse var der fra starten en del problematik, bl.a. var der på forhånd ikke udlagt hvorledes dette skulle foregå. På den baggrund startede forløbet med at der skulle bruges tre uger på indsamling af information omkring hvilken form for teknologi der skulle benyttes. Valget endte med at stå mellem 3d laserskanning eller computer vision.

Det viste sig dog hurtigt at laserskanning var udelukket grundet de høje omkostninger det ville indebære, og der blev derfor fundet et ældre kamera frem der kunne benyttes til projektet.



\section{Produktets specifikationer}

Krav for at produktet rent faktisk kan blive en realitet.

\begin{itemize}
    \item Processen må samlet set ikke tage særlig lang tid, gerne under 4 sekunder.
    \item Målingens præcision skal være på $\pm 0.1$ mm
    \item Omkostningerne skal holde sig under 25.000,- DKr.
\end{itemize}

\section{Teknologivalg}
Denne sektion beskriver de forskellige delelementer projektet består af. Det er både med hensyn til hardware men hvilke udviklingsværktøj der blev benyttet. For at kunne danne et komplet billede, er der information medtaget omkring nogle af beslutninger og deres bevæggrunde. 

\subsection{Kamera}
Der er valgt at benytte sig af et kamera, der på baggrund af en laser der kaster en linje, kan måle forskel i højden fra hvor laserens linje kastes på en termoplastisk markering i forhold til, hvor den rammer ved siden af markeringen.
Kameraet er monokromatisk for at udnytte lyset bedre og har en opløsning på 5 megapixel.

\subsection{Laser}
Laseren kaster en linje ned på markeringsområdet med en styrke der gør det er muligt at kunne aflæse dens lokalisation. Forskellen på hvor laseren rammer på striben og hvor den ikke gør udgør den højde der er interessant.

\subsection{OpenCV}
Til denne proces bliver Open Computer Vision \cite{OpenCV} benyttet, hvilket er et gratis computer vision system der gør det muligt at lave behandling på dataen for at lokalisere laserlinjen på de billeder kameraet tager.

\subsection{Begrundelse}
Valget af kamera og laser er baseret på prisen, et moderne kamera med en tilstrækkelig kvalitet koster under 6.000,- DKr, og en linjelaser koster mellem 2.000,- DKr. og 3.000,- DKr. En laser scannings enhed koster over 15.000,- Dkr. hvilket vil sætte for meget pres på budgettet.

Open Computer Vision er gratis, og kan benyttes i ubegrænset omfang, også i kommercielle forbindelser.

\subsection{C++}
Udviklingen har indtil videre foregået i C$++$, da dette potentielt åbner op for at benytte sig af softwaren, eller dele af den, på generisk vis på tværs af diverse platforme, da softwaren er skrevet med udgangspunkt i standard template library.
Softwaren kræver en compiler der understøtter C++14, for en komplet liste over understøttede compilere, se \href{http://en.cppreference.com/w/cpp/compiler_support}{her}.

\subsection{Udviklingsværktøj}
Softwaren er udviklet ved brug af Microsoft Visual Studio, Community Edition. Selve koden skulle være 100\% kompatibel med andre udviklingskæder som f.eks. GCC. Det ville dog kræve at der blev sat et projekt op der sørger for de korrekte imports.

\subsection{PvAPI}
Kameraet bliver kontrolleret ved hjælp af PvAPI. PvAPI er ellers blevet erstattet af det nyere API fra Allied Vision, Vimba. Men dette API kunne ikke bruges til projektet da det fejler med at kunne kontrollere elementære ting som f.eks. eksponering osv.
Dog skal Vimba's drivere til kameraet været installeret på systemet.


\section{Forsøgsopstillingen}
Denne sektion omhandler opstillingen der er blevet benyttet til at lave test målinger med.



\section{Databehandlingen}



\section{Nuværende resultater}

Denne sektion beskriver de resultater, både de endelige samt de metoder og algoritmer brugt og udviklet under forløbet.

\subsection{}



\section{Videre arbejde}


%\bibliographystyle{unsrt}
%\bibliography{ref}

\end{document}
